\documentclass[Main.tex]{subfiles}
\begin{document}
%%%Phần mở đầu
\chapter{Sinh học cơ thể người}
\section{Dinh dưỡng và tiêu hóa}
\begin{Muctieu}
	\begin{itemize}
		\item Trình bày được khái niệm chất dinh dưỡng và dinh dưỡng.
		\item Kể tên và nêu vai trò của các nhóm chất dinh dưỡng chính.
		\item Mô tả được quá trình tiêu hóa thức ăn ở người.
		\item Nêu được một số bệnh về đường tiêu hóa và cách phòng ngừa.
	\end{itemize}
\end{Muctieu}
\begin{kd}
	\immini{Tại sao chúng ta cần ăn uống? Điều gì xảy ra với thức ăn sau khi chúng ta ăn?}{Ảnh minh họa quá trình ăn uống và tiêu hóa thức ăn}
\end{kd}

**Phần lý thuyết**:
\subsection{Nội dung bài học}
\subsubsection{Chất Dinh Dưỡng và Dinh Dưỡng}
	\Noibat[\maunhan][][][]{ \textbf{Chất dinh dưỡng:} }
		\begin{ghinho}
			\textbf{Chất dinh dưỡng:} là các chất có trong thức ăn, cung cấp nguyên liệu cấu tạo cơ thể và năng lượng cho các hoạt động sống.\\
			Ví dụ: Carbohydrate, chất béo (lipid), protein, vitamin và khoáng chất.
			Lưu ý: Chất nào không cấu tạo cơ thể và không cung cấp năng lượng thì không phải chất dinh dưỡng.
		\end{ghinho}
	\Noibat[\maunhan][][\faStar][]{Các nhóm chất dinh dưỡng chính:}
    \begin{itemize}
        \item \textbf{Carbohydrate:} Nguồn năng lượng chính.
        \item \textbf{Chất béo (Lipid):} Có nhiều trong sữa, phô mai.
        \item \textbf{Protein:} Có nhiều trong thịt, cá, trứng, sữa.
        \item \textbf{Vitamin và khoáng chất:} Có nhiều trong rau, củ, quả.
    \end{itemize}
	\Noibat[\maunhan][][][]{Tại sao cần cung cấp đầy đủ chất dinh dưỡng?}
    \begin{itemize}
        \item Để có đủ nguyên liệu cấu tạo cơ thể.
        \item Để có đủ năng lượng cho các hoạt động sống.
    \end{itemize}
	\Noibat[\maunhan][][][]{\textbf{Dinh dưỡng:}}
    \begin{tomtat}
    	\textbf{Dinh dưỡng} là một quá trình bao gồm thu nhận, biến đổi và sử dụng chất dinh dưỡng để duy trì sự sống của cơ thể.
        \begin{itemize}
            \item Giai đoạn thu nhận: Ăn uống.
            \item Giai đoạn biến đổi: Nhờ hệ tiêu hóa.
            \item Giai đoạn sử dụng: Cơ thể sử dụng chất dinh dưỡng để duy trì sự sống.
        \end{itemize}
    \end{tomtat}

\subsubsection{Tiêu Hóa ở Người}
	\Noibat[\maunhan][][][]{Quá trình tiêu hóa gồm 3 giai đoạn chính:}
    \begin{itemize}
        \item Lấy thức ăn vào.
        \item Biến đổi (tiêu hóa) thức ăn.
        \item Thải chất bã ra.
    \end{itemize}
	\Noibat[\maunhan][][][]{Cấu tạo của hệ tiêu hóa:}
    \begin{itemize}
        \item \textbf{Ống tiêu hóa:}
            \begin{itemize}
                \item Đường ống xuất phát từ miệng và kết thúc ở hậu môn, nơi thức ăn di chuyển qua.
                \item Các bộ phận: Miệng, hầu, thực quản, dạ dày, ruột non, ruột già, hậu môn.
            \end{itemize}
        \item \textbf{Tuyến tiêu hóa:}
            \begin{itemize}
                \item Các cơ quan tiết ra dịch tiêu hóa để hỗ trợ quá trình tiêu hóa thức ăn. Thức ăn không đi qua các tuyến này.
                \item Các bộ phận: Tuyến nước bọt, gan, túi mật, tuyến tụy.
            \end{itemize}
    \end{itemize}
	\Noibat[\maunhan][][][]{Chức năng chính của hệ tiêu hóa:}
    \begin{itemize}
        \item Biến đổi thức ăn thành chất dinh dưỡng cơ thể hấp thụ được.
        \item Loại chất thải ra khỏi cơ thể.
    \end{itemize}
	\Noibat[\maunhan][][][]{Quá trình tiêu hóa thức ăn ở người:}
    \begin{itemize}
        \item \textbf{Tiêu hóa ở khoang miệng:}
            \begin{itemize}
                \item Tiêu hóa cơ học: Răng nhai nghiền nát thức ăn, lưỡi đảo trộn thức ăn.
                \item Tiêu hóa hóa học: Enzyme amylase trong nước bọt phân giải một phần tinh bột chín thành đường maltose.
            \end{itemize}
        \item \textbf{Tiêu hóa ở dạ dày:}
            \begin{itemize}
                \item Tiêu hóa cơ học: Dạ dày co bóp nhào trộn thức ăn.
                \item Tiêu hóa hóa học: Enzyme pepsin phân giải protein thành các đoạn protein ngắn.
                \item Dịch vị: Gồm hydrochloric acid (HCl) diệt khuẩn và enzyme pepsin.
                \item Hydrochloric acid: Giúp tiêu diệt vi khuẩn trong thức ăn.
            \end{itemize}
        \item \textbf{Tiêu hóa ở ruột non:}
            \begin{itemize}
                \item Tiêu hóa hóa học:
                    \begin{itemize}
                        \item Dịch tụy, dịch mật và dịch ruột biến đổi chất dinh dưỡng thành chất đơn giản.
                        \item Carbohydrate $\rightarrow$ glucose, protein $\rightarrow$ amino acid.
                    \end{itemize}
                \item Hấp thụ: Các chất đơn giản được hấp thụ vào máu và mạch bạch huyết qua lông ruột.
            \end{itemize}
        \item \textbf{Tiêu hóa ở ruột già và trực tràng:}
            \begin{itemize}
                \item Hấp thụ: Hấp thụ nước và một số chất dinh dưỡng còn lại.
                \item Tiêu hóa: Không có enzyme tiêu hóa, mà nhờ vi khuẩn phân hủy chất còn lại và lên men tạo thành phân.
                \item Bài tiết: Thải phân ra ngoài qua hậu môn.
            \end{itemize}
    \end{itemize}

\subsubsection{Một Số Bệnh Về Đường Tiêu Hóa}
	\Noibat[\maunhan][][][]{\textbf{Sâu răng}}
    \begin{tomtat}
        \begin{itemize}
        	\item Tình trạng tổn thương mô cứng của răng do vi khuẩn gây ra, tạo thành lỗ nhỏ trên răng.
            \item Nguyên nhân: Do vi khuẩn trong miệng tạo ra acid làm hỏng men răng.
            \item Phòng ngừa:
                \begin{itemize}
                    \item Vệ sinh răng miệng sạch sẽ sau khi ăn và hàng ngày (ít nhất 2 lần).
                    \item Đánh răng đúng cách (ít nhất 2 phút), sử dụng kem đánh răng có flour.
                    \item Hạn chế ăn đồ ngọt.
                \end{itemize}
        \end{itemize}
    \end{tomtat}
	\Noibat[\maunhan][][][]{\textbf{Viêm loét dạ dày tá tràng (Đau dạ dày):} }
    \begin{tomtat}
    	Tổn thương viêm và loét lớp niêm mạc dạ dày hoặc tá tràng.
        \begin{itemize}
            \item Nguyên nhân:
                \begin{itemize}
                    \item Vi khuẩn Helicobacter pylori.
                    \item Sử dụng đồ uống có cồn.
                    \item Ăn uống và sinh hoạt không điều độ.
                \end{itemize}
            \item Triệu chứng:
                \begin{itemize}
                    \item Đau vùng bụng trên rốn.
                    \item Đầy bụng, khó tiêu.
                    \item Ợ chua, ợ nóng.
                \end{itemize}
            \item Phòng ngừa:
                \begin{itemize}
                    \item Ăn chín, uống sôi.
                    \item Không hút thuốc, hạn chế rượu bia.
                    \item Ăn uống điều độ, đúng giờ.
                    \item Giữ tinh thần thoải mái, tránh căng thẳng.
                \end{itemize}
        \end{itemize}
    \end{tomtat}
\subsection{Bài tập}
\begin{dang}{Câu hỏi ôn tập}
\end{dang}
\Noibat[\maunhan][][\faBookmark][]{Ví dụ mẫu}
\begin{vd}
	Chất dinh dưỡng là gì? Kể tên các nhóm chất dinh dưỡng chính và vai trò của chúng.
	\loigiai{
        Chất dinh dưỡng là các chất có trong thức ăn, cung cấp nguyên liệu cấu tạo cơ thể và năng lượng cho các hoạt động sống.
        Các nhóm chất dinh dưỡng chính:
        \begin{itemize}
            \item Carbohydrate: Nguồn năng lượng chính.
            \item Chất béo (Lipid): Dự trữ năng lượng, cấu tạo tế bào.
            \item Protein: Cấu tạo cơ thể, enzyme, hormone.
            \item Vitamin và khoáng chất: Điều hòa hoạt động cơ thể.
        \end{itemize}
	}
\end{vd}

%%%%%=====================Bài tập tự luyện Dạng 1==========================%%%
\Noibat[\maunhan][][\faBook][]{Bài tập tự luyện}

\phan{Bài tập tự luận và trả lời ngắn}
%%%=============SOẠN BT===============%%%
\Opensolutionfile{ansbth}[Ans/LGBT-C05B01_Dang1]
\Opensolutionfile{ansbt}[Ans/AnsBT-C05B01_Dang1]
	\begin{bt}
		Mô tả quá trình tiêu hóa thức ăn ở dạ dày.
		\loigiai{
            Ở dạ dày, thức ăn được tiêu hóa cơ học nhờ các cơn co bóp của dạ dày, nhào trộn thức ăn với dịch vị. Tiêu hóa hóa học nhờ enzyme pepsin phân giải protein thành các đoạn protein ngắn. Dịch vị gồm hydrochloric acid (HCl) diệt khuẩn và enzyme pepsin.
		}
	\end{bt}
    \begin{bt}
		Nêu các biện pháp phòng ngừa bệnh sâu răng.
		\loigiai{
            Các biện pháp phòng ngừa bệnh sâu răng:
            \begin{itemize}
                \item Vệ sinh răng miệng sạch sẽ sau khi ăn và hàng ngày (ít nhất 2 lần).
                \item Đánh răng đúng cách (ít nhất 2 phút), sử dụng kem đánh răng có flour.
                \item Hạn chế ăn đồ ngọt.
            \end{itemize}
		}
	\end{bt}
\Closesolutionfile{ansbt}
\Closesolutionfile{ansbth}
%\bangdapanSA{AnsBT-C05B01_Dang1}
%%%%============Phần trắc nghiệm============%%%
\phan{Trắc nghiệm nhiều lựa chọn}
%%%=============SOẠN EX===============%%%
\Opensolutionfile{ansex}[Ans/LGEX-C05B01_Dang1]
\Opensolutionfile{ans}[Ans/Ans-C05B01_Dang1]
	\begin{ex}
		Chất nào sau đây là nguồn năng lượng chính cho cơ thể?
		\choice
		{Protein}
		{Lipid}
		{\True Carbohydrate}
		{Vitamin}
		\loigiai{Carbohydrate là nguồn năng lượng chính cho cơ thể.}
	\end{ex}
    \begin{ex}
		Enzyme nào sau đây có vai trò phân giải protein ở dạ dày?
		\choice
		{Amylase}
		{Lipase}
		{\True Pepsin}
		{Maltase}
		\loigiai{Enzyme pepsin có vai trò phân giải protein ở dạ dày.}
	\end{ex}
\Closesolutionfile{ans}
\Closesolutionfile{ansex}
%\bangdapan{Ans-C05B01_Dang1}

%%%%%%%%%%%%%%%Trắc nghiệm đúng sai%%%%%%%%%%%%%%%%%%%%%%%%
\phan{Bài tập trắc nghiệm Đúng Sai}
%%%=============SOẠN EXTF===============%%%
\Opensolutionfile{ansex}[Ans/LGTF-C05B01_Dang1]
\Opensolutionfile{ansbook}[Ansbook/AnsTF-C05B01_Dang1]
\Opensolutionfile{ans}[Ans/Tempt-C05B01_Dang1]
	\begin{ex}%[0C1Y1-1]
		Vitamin và khoáng chất cung cấp năng lượng cho cơ thể.
		\choiceTF
		{\True Sai}
		{Đúng}
		{Đúng}
		{Đúng}
		\loigiai{
			\begin{itemchoice}[T1,F2,F3,F4]
				\itemch Vitamin và khoáng chất không cung cấp năng lượng, mà có vai trò điều hòa hoạt động của cơ thể.
				\itemch Vitamin và khoáng chất không cung cấp năng lượng, mà có vai trò điều hòa hoạt động của cơ thể.
				\itemch Vitamin và khoáng chất không cung cấp năng lượng, mà có vai trò điều hòa hoạt động của cơ thể.
				\itemch Vitamin và khoáng chất không cung cấp năng lượng, mà có vai trò điều hòa hoạt động của cơ thể.
			\end{itemchoice}
		}
	\end{ex}
    \begin{ex}%[0C1Y1-1]
		Ruột già là nơi hấp thụ chất dinh dưỡng chính của cơ thể.
		\choiceTF
		{\True Sai}
		{Đúng}
		{Đúng}
		{Đúng}
		\loigiai{
			\begin{itemchoice}[T1,F2,F3,F4]
				\itemch Ruột non là nơi hấp thụ chất dinh dưỡng chính của cơ thể. Ruột già chủ yếu hấp thụ nước và một số chất còn lại.
				\itemch Ruột non là nơi hấp thụ chất dinh dưỡng chính của cơ thể. Ruột già chủ yếu hấp thụ nước và một số chất còn lại.
				\itemch Ruột non là nơi hấp thụ chất dinh dưỡng chính của cơ thể. Ruột già chủ yếu hấp thụ nước và một số chất còn lại.
				\itemch Ruột non là nơi hấp thụ chất dinh dưỡng chính của cơ thể. Ruột già chủ yếu hấp thụ nước và một số chất còn lại.
			\end{itemchoice}
		}
	\end{ex}
\Closesolutionfile{ans}
\Closesolutionfile{ansbook}
\Closesolutionfile{ansex}
%\bangdapanTF{AnsTF-C05B01_Dang1}

%%%%%============BTS_1================%%%%%%
\phan{Bài tập trả lời ngắn}
%%%=============SOẠN BT===============%%%
\Opensolutionfile{ansbth}[Ans/LGBT-C05B01_Dang1]
\Opensolutionfile{ansbt}[Ans/AnsBT-C05B01_Dang1]
    \begin{bt}
        Kể tên các tuyến tiêu hóa chính và cho biết vai trò của chúng.
        \shortans{Tuyến nước bọt, gan, túi mật, tuyến tụy}
        \loigiai{
            Các tuyến tiêu hóa chính:
            \begin{itemize}
                \item Tuyến nước bọt: Tiết nước bọt chứa enzyme amylase phân giải tinh bột.
                \item Gan: Sản xuất dịch mật giúp tiêu hóa chất béo.
                \item Túi mật: Lưu trữ và cô đặc dịch mật.
                \item Tuyến tụy: Tiết dịch tụy chứa các enzyme tiêu hóa carbohydrate, protein và lipid.
            \end{itemize}
        }
    \end{bt}
\Closesolutionfile{ansbt}
\Closesolutionfile{ansbth}
\end{document}